\SourceLeaf{19}
\begin{center}
{\large\bfseries ZWEITER ABSCHNITT}\par\bigskip
{\small\bfseries ÜBER}\par\bigskip
{\large\bfseries KONGRUENZEN ERSTEN GRADES.}
\end{center}
\begin{center}\rule{0.09\linewidth}{0.3pt}\end{center}
\begin{center}\emph{Vorbereitende Sätze über Primzahlen, Faktoren usw.}\end{center}
\Article{13}
\textsc{Satz}. \emph{Das Produkt zweier positiver Zahlen, die kleiner als eine gegebene Primzahl sind, kann nicht durch diese Primzahl geteilt werden.}

Sei $p$ eine Primzahl und $a$ positiv und $<p$: Dann gibt es keine positive Zahl $b$, die kleiner als $p$ ist, sodass $ab\equiv 0\ (\mathrm{mod.}\ p)$ ist.

\emph{Beweis.} Wenn jemand dies bestreitet, nehmen wir an, es gebe Zahlen $b, c, d$ usw., alle $<p$, für die $ab\equiv 0$, $ac\equiv 0$, $ad\equiv 0$ usw. (mod. $p$) ist. Sei $b$ die kleinste von allen, sodass keine Zahl, die kleiner als $b$ ist, diese Eigenschaft hat. Offensichtlich ist $b>1$: Wäre nämlich $b=1$, so wäre $ab=a<p$ (\emph{nach Voraussetzung}) und damit nicht durch $p$ teilbar. Da $p$ eine Primzahl ist, kann es also nicht durch $b$ geteilt werden, sondern liegt zwischen den beiden benachbarten Vielfachen von $b$, $mb$ und $(m+1)b$. Sei $p-mb=b'$; dann ist $b'$ eine positive Zahl und $<b$. Da wir nun $ab\equiv 0\ (\mathrm{mod.}\ p)$ vorausgesetzt haben, gilt auch $mab\equiv 0$ (Art. 7); zieht man dies von $ap\equiv 0$ ab, so erhält man $a(p-mb)=ab'\equiv 0$. Also muss $b'$ zu den Zahlen $b, c, d$ usw. gehören, obwohl es kleiner als ihre kleinste, $b$, ist. \emph{Q.~E.~A.}

\Article{14}
\emph{Kann weder $a$ noch $b$ durch die Primzahl $p$ geteilt werden, so kann auch das Produkt $ab$ nicht durch $p$ geteilt werden.}

\SourceLeaf{20}
Seien die kleinsten positiven Reste der Zahlen $a, b$ nach dem Modul $p$ mit $\alpha, \beta$ bezeichnet; keiner von ihnen ist $0$ (\emph{nach Voraussetzung}). Wäre nun $ab\equiv 0$ (mod. $p$), so wäre wegen $ab\equiv\alpha\beta$ auch $\alpha\beta\equiv 0$, was mit dem vorigen Satz unvereinbar ist.

Einen Beweis dieses Satzes hat bereits Euklid gegeben, \emph{El.} VII. 32. Wir wollten ihn trotzdem nicht weglassen: zum einen, weil mehrere Neuere entweder vage Überlegungen als Beweis ausgegeben oder den Satz ganz übergangen haben; zum anderen, weil sich die Art der hier verwendeten Methode, die wir später zur Klärung viel verborgenerer Dinge benutzen werden, an einem einfacheren Fall leichter verstehen lässt.

\Article{15}
\emph{Kann keine der Zahlen $a, b, c, d$ usw. durch die Primzahl $p$ geteilt werden, so kann auch das Produkt $abcd$ usw. nicht durch $p$ geteilt werden.}

Nach dem vorigen Artikel ist $ab$ nicht durch $p$ teilbar; also auch nicht $abc$ und daher auch nicht $abcd$ usw.

\Article{16}
\textsc{Satz}. \emph{Jede zusammengesetzte Zahl lässt sich nur auf eine einzige Weise in Primfaktoren zerlegen.}

\emph{Beweis.} Dass sich jede zusammengesetzte Zahl in Primfaktoren zerlegen lässt, ist aus den Anfangsgründen bekannt; dass dies aber nicht auf mehrere verschiedene Weisen geschehen kann, wird meist zu Unrecht stillschweigend vorausgesetzt. Nehmen wir an, eine zusammengesetzte Zahl $A$, die $=a^\alpha b^\beta c^\gamma$ usw. ist, wobei $a, b, c$ usw. verschiedene Primzahlen bezeichnen, ließe sich noch auf eine andere Weise in Primfaktoren zerlegen. Zunächst ist klar, dass in diesem zweiten System von Faktoren keine anderen Primzahlen als $a, b, c$ usw. vorkommen können, weil keine andere Primzahl die aus diesen zusammengesetzte Zahl $A$ teilen kann. Ebenso darf in diesem zweiten System keine der Primzahlen $a, b, c$ usw. fehlen, da sie sonst $A$ nicht teilen würde (voriger Art.). Die beiden Zerlegungen können sich also nur dadurch unterscheiden, dass irgendeine Primzahl in der einen öfter vorkommt als in der anderen. Sei $p$ eine solche Primzahl, die in der einen Zerlegung $m$-mal und in der anderen $n$-mal vorkommt, und sei $m>n$: Nun werde der Faktor $p$ aus beiden Systemen $n$-mal entfernt, sodass er in dem einen noch $m-n$-mal vorkommt, aus dem anderen aber ganz verschwunden ist. Es gibt also zwei Zerlegungen der Zahl $\frac{A}{p^n}$ in Faktoren, von denen die eine den Faktor $p$ gar nicht, die andere aber $m-n$-mal enthält, entgegen dem eben Bewiesenen.

\SourceLeaf{21}
\Article{17}
Ist also eine zusammengesetzte Zahl $A$ das Produkt von $B, C, D$ usw., so können unter den Primfaktoren der Zahlen $B, C, D$ usw. offensichtlich keine anderen vorkommen als unter den Faktoren von $A$; und jeder dieser Faktoren muss in $B, C, D$ usw. zusammengenommen ebenso oft vorkommen wie in $A$. Daraus ergibt sich ein Kriterium dafür, ob eine Zahl $B$ eine andere Zahl $A$ teilt oder nicht. Sie tut dies, wenn $B$ weder andere Primfaktoren noch einen von ihnen öfter als $A$ enthält; ist eine dieser Bedingungen nicht erfüllt, so teilt $B$ die Zahl $A$ nicht.

Ist $A=a^\alpha b^\beta c^\gamma$ usw., wobei $a, b, c$ usw. wie oben verschiedene Primzahlen bezeichnen, so lässt sich daraus durch Abzählen der Kombinationen leicht ableiten, dass~$A$
\[
(\alpha+1)(\beta+1)(\gamma+1)\ \mathrm{etc.}
\]
verschiedene Teiler besitzt, einschließlich $1$ und $A$.

\Article{18}
Ist also $A=a^\alpha b^\beta c^\gamma$ usw., $K=k^\kappa l^\lambda m^{\mu}$ usw., und sind die Primzahlen $a, b, c$ usw., $k, l, m$ usw. alle verschieden, so haben $A$ und $K$ offensichtlich außer $1$ keinen gemeinsamen Teiler, sind also teilerfremd.

Sind mehrere Zahlen $A, B, C$ usw. gegeben, so wird ihr \emph{größter gemeinsamer Teiler} folgendermaßen bestimmt. Man zerlegt sie alle in ihre Primfaktoren und wählt darunter diejenigen aus, die allen Zahlen $A, B, C$ usw. gemeinsam sind (gibt es solche nicht, so gibt es keinen allen gemeinsamen Teiler). Dann notiert man, wie oft jeder dieser Primfaktoren in jeder der Zahlen $A, B, C$ usw. enthalten ist, also \emph{welchen Exponenten} jeder in den einzelnen Zahlen $A, B, C$ usw. hat. Schließlich gibt man jedem Primfaktor den kleinsten seiner Exponenten in $A, B, C$ usw. und bildet daraus das Produkt; dies ist der gesuchte gemeinsame Teiler.

Sucht man dagegen das \emph{kleinste gemeinsame Vielfache} der Zahlen $A, B, C$ usw., so geht man folgendermaßen vor. Man sammelt alle Primzahlen, die wenigstens eine der Zahlen $A, B, C$ usw. teilen, gibt jeder den größten ihrer Exponenten in den Zahlen $A, B, C$ usw. und bildet so das Produkt aus allen; dies ist das gesuchte Vielfache.

\emph{Beispiel.} Sei $A=504=2^3 3^2 7$; $B=2880=2^6 3^2 5$; $C=864=2^5 3^3$. Für den größten gemeinsamen Teiler hat man die Primfaktoren $2, 3$, denen die Exponenten $3, 2$ zu geben sind; es ergibt sich $=2^3 3^2=72$; das kleinste gemeinsame Vielfache ist dagegen $2^6 3^3 5.7=60480$.

\SourceLeaf{22}
Wegen ihrer Einfachheit lassen wir die Beweise weg. Wie diese Aufgaben zu lösen sind, wenn die Zerlegung der Zahlen $A, B, C$ usw. in Faktoren nicht gegeben ist, ist im Übrigen aus den Anfangsgründen bekannt.

\Article{19}
\emph{Sind die Zahlen $a, b, c$ usw. zu einer anderen Zahl $k$ teilerfremd, so ist auch ihr Produkt $abc$ usw. zu $k$ teilerfremd.}

Da nämlich keine der Zahlen $a, b, c$ usw. einen Primfaktor mit $k$ gemeinsam hat und das Produkt $abc$ usw. keine anderen Primfaktoren haben kann als die Faktoren einer der Zahlen $a, b, c$ usw., hat auch das Produkt $abc$ usw. keinen Primfaktor mit $k$ gemeinsam. Nach dem vorigen Artikel ist also $k$ zu $abc$ usw. teilerfremd.

\emph{Sind die Zahlen $a, b, c$ usw. paarweise teilerfremd und teilt jede eine andere Zahl $k$, so teilt auch ihr Produkt die Zahl $k$.}

Dies folgt ebenso leicht aus den Artikeln 17, 18. Sei nämlich ein beliebiger Primteiler des Produkts $abc$ usw. mit $p$ bezeichnet, der darin $\pi$-mal vorkomme; offensichtlich muss eine der Zahlen $a, b, c$ usw. denselben Teiler $\pi$-mal enthalten. Daher enthält auch $k$, das von dieser Zahl geteilt wird, $\pi$-mal den Teiler $p$. Ebenso verhält es sich mit den übrigen Teilern des Produkts $abc$ usw.

Daher gilt: \emph{Sind zwei Zahlen $m, n$ nach mehreren paarweise teilerfremden Moduln $a, b, c$ usw. kongruent, so sind sie auch nach deren Produkt kongruent.} Da nämlich $m-n$ durch jede der Zahlen $a, b, c$ usw. teilbar ist, lässt es sich auch durch ihr Produkt teilen.

Schließlich gilt: \emph{Ist $a$ zu $b$ teilerfremd und $ak$ durch $b$ teilbar, so ist auch $ak$ durch $b$ teilbar.} Denn da $ak$ sowohl durch $a$ als auch durch $b$ teilbar ist, lässt es sich auch durch $ab$ teilen; das heißt, $\frac{ak}{ab}=\frac{k}{b}$ ist ganzzahlig.

\Article{20}
\emph{Ist $A=a^\alpha b^\beta c^\gamma$ usw., wobei $a, b, c$ usw. verschiedene Primzahlen bezeichnen, irgendeine Potenz, etwa $=k^n$, so sind alle Exponenten $\alpha, \beta, \gamma$ usw. durch $n$ teilbar.}

Die Zahl $k$ enthält nämlich keine anderen Primfaktoren als $a, b, c$ usw. Enthält sie den Faktor $a$ genau $\alpha'$-mal, so enthält $k^n$, also $A$, diesen Faktor $n\alpha'$-mal; daher ist $n\alpha'=\alpha$ und $\frac{\alpha}{n}$ ganzzahlig. Ebenso zeigt man, dass $\frac{\beta}{n}$ usw. ganzzahlig sind.
